\documentclass[11pt]{article}
\usepackage[margin=1in]{geometry}
\usepackage{amsmath,amssymb,amsthm}
\usepackage{booktabs}
\usepackage{hyperref}

\newtheorem{theorem}{Theorem}
\newtheorem{proposition}{Proposition}
\newtheorem{corollary}{Corollary}
\newtheorem{lemma}{Lemma}
\theoremstyle{remark}
\newtheorem*{remark}{Remark}
\theoremstyle{definition}
\newtheorem*{example}{Example}

\title{The Brock-Hommes Adaptive Belief System: A Complete Mathematical Argument}
\author{Wulin Suo}
\date{}

\begin{document}
\maketitle

\begin{abstract}
\noindent This note derives the Brock-Hommes adaptive belief system from first principles: the mean-variance (CARA) demand of a heterogeneous population of investors, the market-clearing price equation their demands imply, the recasting of that equation in deviation-from-fundamental form, and the discrete-choice (logit) rule by which the population's belief mix evolves in response to each strategy's realized performance. Specializing to the canonical two-type case, a costly fundamentalist rule against a free trend-extrapolating rule, it derives, in closed form, the equilibrium population split at the fundamental steady state and the exact condition under which that steady state is locally stable. The central result is a genuine bifurcation in the model's own intensity-of-choice parameter: increasing agents' sensitivity to past performance differences can destabilize the very steady state that made the two rules equally accurate in the first place, the mechanism Brock and Hommes (1997) name a ``rational route to randomness.'' The note closes with a worked numerical example, a discussion of the model's assumptions and extensions, and its relationship to the simpler, fixed-population fundamentalist-chartist example used elsewhere in this book.
\end{abstract}

\section{The Problem: Where Do Trading-Population Shares Come From?}
\label{sec:problem}

A heterogeneous-agent asset-pricing model needs some population of fundamentalists and some population of trend-chasers to produce anything interesting; the natural next question is where those population shares come from, and whether they should be treated as fixed parameters at all. Real investors are not assigned a trading style at birth: a strategy that has recently performed well attracts new adherents, and one that has recently performed poorly loses them. Brock and Hommes (1997, 1998) formalize this as an \emph{evolutionary} or \emph{adaptive belief system}: a population of investors chooses, each period, among a finite menu of forecasting rules using a discrete-choice model driven by each rule's own recently realized fitness, so the population weights $n_{h,t}$ that a fixed-mix model treats as exogenous become, in this model, an endogenous output of rational, performance-based rule selection. This note derives that system in full and shows that the resulting feedback, prices depend on population weights, population weights depend on past prices, produces a genuine bifurcation: as agents become more sensitive to performance differences (formally, as an ``intensity of choice'' parameter $\beta$ rises), a market that started out stable can lose stability entirely on its own, without any change to the underlying trading rules themselves.

\section{The Asset-Pricing Model with Heterogeneous Beliefs}
\label{sec:model}

\subsection{Assumptions}

\begin{itemize}
\item[(A1)] A risky asset trades at ex-dividend price $p_t$ and pays a stochastic dividend $y_t$ at each date $t$; a risk-free asset offers a constant gross return $R=1+r>1$ per period, in perfectly elastic supply.
\item[(A2)] Each investor $i$ has constant absolute risk aversion (CARA) utility $U(W)=-e^{-aW}$, $a>0$, and chooses a number of risky shares $z_t$ at date $t$ to maximize $E_{i,t}[U(W_{t+1})]$, where wealth evolves as $W_{t+1}=RW_t+z_t D_{t+1}$, $D_{t+1}\equiv p_{t+1}+y_{t+1}-Rp_t$ the excess (over the risk-free alternative) payoff per share.
\item[(A3)] Investors fall into finitely many belief types $h=1,\dots,H$. Every investor of type $h$ believes $D_{t+1}$ is, conditional on $\mathcal F_t$, normally distributed with a common variance $\sigma^2$ (the same across types) and a type-specific conditional mean $d_{h,t}\equiv E_{h,t}[p_{t+1}+y_{t+1}]-Rp_t$.
\item[(A4)] The risky asset is in zero net supply: at every date, the fractions $n_{h,t}\ge0$, $\sum_h n_{h,t}=1$, of investors holding each belief type must clear the market, $\sum_h n_{h,t}z_{h,t}=0$.
\end{itemize}

\begin{lemma}[CARA demand under conditional normality]
\label{lem:cara}
Under (A2)--(A3), type $h$'s optimal date-$t$ demand is
\begin{equation}
z_{h,t} = \frac{d_{h,t}}{a\sigma^2} = \frac{E_{h,t}[p_{t+1}+y_{t+1}]-Rp_t}{a\sigma^2}.
\label{eq:demand}
\end{equation}
\end{lemma}
\begin{proof}
Fix a candidate demand $z$. Since $D_{t+1}\mid\mathcal F_t \sim_h N(d_{h,t},\sigma^2)$ and $W_{t+1}=RW_t+zD_{t+1}$,
\begin{equation*}
E_{h,t}\big[-e^{-aW_{t+1}}\big] = -e^{-aRW_t}\,E_{h,t}\big[e^{-azD_{t+1}}\big] = -e^{-aRW_t}\exp\!\Big(-az\,d_{h,t}+\tfrac12 a^2z^2\sigma^2\Big),
\end{equation*}
using the normal moment-generating function $E[e^{sX}]=\exp(s\mu+\tfrac12s^2\sigma^2)$ for $X\sim N(\mu,\sigma^2)$ with $s=-az$. Since $u\mapsto -e^{-u}$ is strictly increasing, maximizing $E_{h,t}[-e^{-aW_{t+1}}]$ over $z$ is equivalent to maximizing $azd_{h,t}-\tfrac12a^2z^2\sigma^2$, equivalently (dividing by $a>0$) maximizing $zd_{h,t}-\tfrac12az^2\sigma^2$, a strictly concave quadratic in $z$ with first-order condition $d_{h,t}-az\sigma^2=0$, giving the unique maximizer $z_{h,t}=d_{h,t}/(a\sigma^2)$.
\end{proof}

\begin{proposition}[Market-clearing price equation]
\label{prop:clearing}
Under (A1)--(A4),
\begin{equation}
Rp_t = \sum_{h=1}^{H} n_{h,t}\,E_{h,t}[p_{t+1}+y_{t+1}].
\label{eq:clearing}
\end{equation}
\end{proposition}
\begin{proof}
Substituting Equation~\eqref{eq:demand} into the market-clearing condition of (A4),
\begin{equation*}
0 = \sum_h n_{h,t}z_{h,t} = \frac{1}{a\sigma^2}\left(\sum_h n_{h,t}E_{h,t}[p_{t+1}+y_{t+1}] - Rp_t\sum_h n_{h,t}\right).
\end{equation*}
Since $\sum_h n_{h,t}=1$, this is equivalent to Equation~\eqref{eq:clearing}.
\end{proof}

Equation~\eqref{eq:clearing} says the current price is the (population-weighted) average of every type's forecast of tomorrow's payoff, discounted at the risk-free rate, exactly the risk-neutral pricing equation that would hold under a single, representative rational agent, except that the forecast being discounted is now a mixture across genuinely different beliefs rather than a single correct one.

\section{The Fundamental Benchmark and Price Deviations}
\label{sec:fundamental}

\begin{itemize}
\item[(A5)] All types agree on the dividend process and on a common ``fundamental'' price process $p_t^*$ satisfying, for every $t$, $Rp_t^*=E_t[p_{t+1}^*+y_{t+1}]$ together with the transversality (no-bubble) condition $\lim_{k\to\infty}R^{-k}E_t[p_{t+k}^*]=0$.
\item[(A6)] Type $h$'s belief decomposes as $E_{h,t}[p_{t+1}+y_{t+1}] = E_t[p_{t+1}^*+y_{t+1}] + f_h(x_{t-1})$, where $x_t \equiv p_t-p_t^*$ is the price's deviation from fundamental and $f_h$ is type $h$'s (deterministic) forecasting rule, applied to the most recently observed deviation.
\end{itemize}

\begin{lemma}[The fundamental price is a discounted dividend stream]
\label{lem:fundamental}
Under (A5), $p_t^* = \sum_{k=1}^{\infty} R^{-k}E_t[y_{t+k}]$.
\end{lemma}
\begin{proof}
By induction on $K$, $p_t^*=\sum_{k=1}^{K}R^{-k}E_t[y_{t+k}]+R^{-K}E_t[p_{t+K}^*]$. The base case $K=1$ is (A5)'s recursion divided by $R$. For the inductive step, apply (A5) at date $t+K$ and the law of iterated expectations, $E_t[p_{t+K}^*]=E_t\big[R^{-1}E_{t+K}[p_{t+K+1}^*]+R^{-1}E_{t+K}[y_{t+K+1}]\big]=R^{-1}E_t[p^*_{t+K+1}]+R^{-1}E_t[y_{t+K+1}]$, which extends the sum to $K+1$ terms. Letting $K\to\infty$ and invoking the transversality condition of (A5) gives the stated result.
\end{proof}

\begin{theorem}[Deviation-form price equation]
\label{thm:deviation}
Under (A1)--(A6),
\begin{equation}
x_t = \frac{1}{R}\sum_{h=1}^{H} n_{h,t}\,f_h(x_{t-1}).
\label{eq:deviation}
\end{equation}
\end{theorem}
\begin{proof}
Substituting (A6) into Equation~\eqref{eq:clearing} and using $\sum_h n_{h,t}=1$,
\begin{equation*}
Rp_t = \sum_h n_{h,t}\Big(E_t[p_{t+1}^*+y_{t+1}]+f_h(x_{t-1})\Big) = E_t[p_{t+1}^*+y_{t+1}] + \sum_h n_{h,t}f_h(x_{t-1}).
\end{equation*}
By (A5), $E_t[p_{t+1}^*+y_{t+1}]=Rp_t^*$, so $Rp_t-Rp_t^* = \sum_h n_{h,t}f_h(x_{t-1})$; dividing by $R$ and using $x_t=p_t-p_t^*$ gives Equation~\eqref{eq:deviation}.
\end{proof}

Equation~\eqref{eq:deviation} is the model's central pricing relation: the current deviation from fundamental is a weighted average of every type's forecast of the deviation, formed from the single most recently observed value, discounted at the risk-free rate. It is a strict generalization of the fixed-weight fundamentalist-chartist recursion used elsewhere in this book: setting $n_{h,t}\equiv n_h$ constant for all $t$ recovers a fixed-mix model of exactly that kind. Section~\ref{sec:switching} makes the weights $n_{h,t}$ themselves an equilibrium object.

\section{Evolutionary Fitness and the Discrete-Choice Switching Rule}
\label{sec:switching}

\begin{itemize}
\item[(A7)] Type $h$ incurs a (possibly zero) per-period cost $C_h\ge0$ for using its forecasting rule, meant to capture the cost of gathering and processing the information a more sophisticated rule requires. Type $h$'s realized fitness at date $t-1$ is the negative squared error of its date-$(t-2)$ forecast of the date-$(t-1)$ deviation, net of cost,
\begin{equation}
U_{h,t-1} = -\big(x_{t-1}-f_h(x_{t-2})\big)^2 - C_h.
\label{eq:fitness}
\end{equation}
\item[(A8)] Population fractions update via a discrete-choice (logit) rule with intensity of choice $\beta\ge0$,
\begin{equation}
n_{h,t} = \frac{\exp(\beta U_{h,t-1})}{\sum_{h'=1}^{H}\exp(\beta U_{h',t-1})}.
\label{eq:logit}
\end{equation}
\end{itemize}

Larger $\beta$ makes the population more sensitive to small fitness differences ($\beta\to\infty$ sends essentially all agents to whichever rule performed even marginally better last period), while $\beta=0$ ignores performance entirely and splits the population evenly regardless of it.

\begin{lemma}[Binary logit]
\label{lem:logit}
For $H=2$, writing $\Delta U_{t-1}\equiv U_{2,t-1}-U_{1,t-1}$,
\begin{equation}
n_{2,t} = \frac{1}{1+e^{-\beta \Delta U_{t-1}}}, \qquad n_{1,t}=1-n_{2,t}.
\label{eq:binarylogit}
\end{equation}
\end{lemma}
\begin{proof}
Dividing numerator and denominator of Equation~\eqref{eq:logit} (for $h=2$) by $\exp(\beta U_{2,t-1})$,
\begin{equation*}
n_{2,t} = \frac{1}{1+\exp\!\big(\beta(U_{1,t-1}-U_{2,t-1})\big)} = \frac{1}{1+e^{-\beta\Delta U_{t-1}}}.
\end{equation*}
$n_{1,t}=1-n_{2,t}$ follows since $H=2$ and the two fractions sum to one.
\end{proof}

\section{The Two-Type Fundamentalist/Trend-Chaser System}
\label{sec:twotype}

Specialize to $H=2$: type $1$ (``fundamentalist'') forecasts no deviation at all, $f_1(x)\equiv0$, and pays cost $C_1=C\ge0$; type $2$ (``trend-chaser'') linearly extrapolates the most recent deviation, $f_2(x)=gx$ for a constant $g$, and pays no cost, $C_2=0$. Substituting into Equation~\eqref{eq:deviation}, and since $f_1\equiv0$ contributes nothing,
\begin{equation}
x_t = \frac{g}{R}\,n_{2,t}\,x_{t-1}.
\label{eq:twotypeprice}
\end{equation}
From Equation~\eqref{eq:fitness}, $U_{1,t-1}=-x_{t-1}^2-C$ and $U_{2,t-1}=-(x_{t-1}-gx_{t-2})^2$, so
\begin{equation}
\Delta U_{t-1} = x_{t-1}^2+C-(x_{t-1}-gx_{t-2})^2.
\label{eq:twotypedelta}
\end{equation}
Equations~\eqref{eq:twotypeprice}--\eqref{eq:twotypedelta} together with Lemma~\ref{lem:logit} define a second-order nonlinear recursion, $x_t=\Phi(x_{t-1},x_{t-2})$, in the single observable $x_t$.

\section{Local Stability of the Fundamental Steady State}
\label{sec:stability}

\begin{lemma}[The fundamental steady state]
\label{lem:steadystate}
$x_t\equiv0$ is a steady state of the system of Section~\ref{sec:twotype} for every $\beta,C,g$, with associated constant population fraction
\begin{equation}
n_2^{*} = \frac{1}{1+e^{-\beta C}}.
\label{eq:nstar}
\end{equation}
\end{lemma}
\begin{proof}
At $x_{t-1}=x_{t-2}=0$, Equation~\eqref{eq:twotypedelta} gives $\Delta U^{*}=0+C-0=C$, so Lemma~\ref{lem:logit} gives $n_2^{*}=1/(1+e^{-\beta C})$, a constant. Equation~\eqref{eq:twotypeprice} then gives $x_t = \tfrac{g}{R}n_2^{*}\cdot 0 = 0$, confirming $x_t\equiv0$ is self-consistent.
\end{proof}

At the fundamental steady state, both rules forecast exactly the same thing (zero deviation), so their forecast-error fitness differs only by the cost $C$; Equation~\eqref{eq:nstar} says the equilibrium population split is exactly the logit choice between two equally accurate, differently priced options.

\begin{theorem}[Local stability]
\label{thm:stability}
The fundamental steady state $x_t\equiv0$ is locally asymptotically stable if $\left|g\,n_2^{*}/R\right|<1$ and unstable if $\left|g\,n_2^{*}/R\right|>1$, where $n_2^{*}$ is given by Equation~\eqref{eq:nstar}.
\end{theorem}
\begin{proof}
Write $\Phi(u,v) = \frac{g}{R}\,u\,n_2(u,v)$ for the right-hand side of Equation~\eqref{eq:twotypeprice}, where $n_2(u,v)=\big[1+\exp(-\beta\Delta U(u,v))\big]^{-1}$ and $\Delta U(u,v) = u^2+C-(u-gv)^2$ is Equation~\eqref{eq:twotypedelta} restated as a function of $(x_{t-1},x_{t-2})=(u,v)$. Both $\Delta U$ and $n_2$ are continuously differentiable in $(u,v)$. Differentiating $\Phi$,
\begin{equation*}
\frac{\partial \Phi}{\partial u}(u,v) = \frac{g}{R}\left[n_2(u,v) + u\,\frac{\partial n_2}{\partial u}(u,v)\right], \qquad
\frac{\partial \Phi}{\partial v}(u,v) = \frac{g}{R}\,u\,\frac{\partial n_2}{\partial v}(u,v).
\end{equation*}
Evaluating at $(u,v)=(0,0)$: the term $u\,\partial n_2/\partial u$ vanishes because $u=0$, regardless of the (finite) value of the derivative it multiplies, leaving $\partial\Phi/\partial u(0,0) = \tfrac{g}{R}n_2(0,0)=\tfrac{g}{R}n_2^{*}$; by the identical argument, $\partial\Phi/\partial v(0,0)=0$. The linearization of the second-order recursion $x_t=\Phi(x_{t-1},x_{t-2})$ around $(0,0)$ is therefore
\begin{equation*}
x_t \approx \frac{g\,n_2^{*}}{R}\,x_{t-1} + 0\cdot x_{t-2},
\end{equation*}
a first-order linear recursion with a single nonzero characteristic root $\lambda=g\,n_2^{*}/R$ (the companion-matrix form of the linearized second-order system has eigenvalues $\lambda$ and $0$). A linear recursion converges to zero from any initial condition if and only if every characteristic root has modulus strictly less than one, which holds here if and only if $|\lambda|<1$; standard linearization theory (the root of the Jacobian of the map at the fixed point determining local stability) then transfers this conclusion to the original nonlinear recursion in a neighborhood of $(0,0)$.
\end{proof}

\begin{corollary}[Three stability regimes]
\label{cor:regimes}
Assume $g>0$ and $C>0$. Since $n_2^{*}\in(1/2,1)$ is strictly increasing in $\beta$ with $n_2^{*}\to \tfrac12$ as $\beta\to0^{+}$ and $n_2^{*}\to1$ as $\beta\to\infty$:
\begin{itemize}
\item[(i)] If $g\le R$, the fundamental steady state is stable for every $\beta\ge0$.
\item[(ii)] If $R<g<2R$, there is a unique critical intensity of choice
\begin{equation}
\beta^{*} = \frac{1}{C}\ln\!\left(\frac{R}{g-R}\right) > 0,
\label{eq:betastar}
\end{equation}
such that the steady state is stable for $\beta<\beta^{*}$ and unstable for $\beta>\beta^{*}$.
\item[(iii)] If $g\ge2R$, the fundamental steady state is unstable for every $\beta\ge0$.
\end{itemize}
\end{corollary}
\begin{proof}
By Theorem~\ref{thm:stability}, stability holds iff $g n_2^{*}<R$ (since $g,n_2^{*},R>0$), i.e., iff $n_2^{*}<R/g$. At $\beta=0$, $n_2^{*}=1/2$; since $C>0$, $\Delta U^{*}=C>0$ for every $\beta>0$, and Equation~\eqref{eq:nstar} is strictly increasing in $\beta$, approaching $1$ as $\beta\to\infty$. If $g\le R$ then $R/g\ge1>n_2^{*}$ for every $\beta$, giving (i). If $g\ge 2R$ then $R/g\le1/2\le n_2^{*}$ for every $\beta\ge0$, giving (iii). If $R<g<2R$ then $1/2<R/g<1$, so by continuity and strict monotonicity of $n_2^{*}(\beta)$ there is a unique $\beta^{*}>0$ solving $n_2^{*}(\beta^{*})=R/g$; solving Equation~\eqref{eq:nstar} for $\beta$ at this value, $1+e^{-\beta^{*}C}=g/R \iff e^{-\beta^{*}C}=(g-R)/R \iff \beta^{*}=\tfrac1C\ln\big(R/(g-R)\big)$, which is positive precisely because $0<(g-R)/R<1$ in this regime. Since $n_2^{*}(\beta)$ is strictly increasing, $n_2^{*}(\beta)<R/g$ (stability) holds exactly for $\beta<\beta^{*}$.
\end{proof}

\begin{remark}[Why this is a ``rational route to randomness'']
Corollary~\ref{cor:regimes}(ii) is the model's central qualitative message. In that regime the market is stable at low intensity of choice, exactly where the population is roughly evenly split and trend-chasing is too diluted to destabilize the price, yet becomes unstable purely because agents grow \emph{more} responsive to (correctly perceived, genuinely existing) performance differences between the two rules, not because the rules themselves change. At the fundamental steady state the two rules are equally accurate, so a more rational, more performance-sensitive population rationally migrates toward the free rule over the costly one, and it is exactly this rational migration that erodes the stabilizing weight of the costly, but otherwise-equivalent, fundamentalist rule. Brock and Hommes (1997) show, beyond the local result proved here, that once $\beta$ crosses $\beta^{*}$ the map's dynamics do not merely diverge from the steady state but undergo a period-doubling bifurcation route to deterministic chaos, i.e., what looks statistically like unpredictable, noisy price behavior is in fact generated by a fully deterministic, rational, low-dimensional system, the ``rational route to randomness'' of the paper's title.
\end{remark}

\begin{remark}[A cost asymmetry is essential]
If $C=0$, Equation~\eqref{eq:nstar} gives $n_2^{*}=1/2$ for every $\beta$, so the stability boundary in Corollary~\ref{cor:regimes} collapses to the single, $\beta$-independent condition $g<2R$: without some genuine cost or quality asymmetry between the two rules, increasing the population's sensitivity to performance has nothing to bite on at the steady state, since both rules are tied there for every $\beta$. The bifurcation of Corollary~\ref{cor:regimes}(ii) requires both a positive cost differential and a sufficiently strong trend-extrapolation parameter $g$; the direction of the effect also requires the cost to fall on the \emph{stabilizing} rule, if instead $C_1=0<C_2=C$ (the trend-chaser is the costly one), the same algebra gives $n_2^{*}\to0$ as $\beta\to\infty$, and rising $\beta$ makes the market \emph{more}, not less, stable.
\end{remark}

\section{Worked Numerical Example}
\label{sec:example}

Let $R=1.01$ (a one percent per-period risk-free rate), $g=1.5$, and $C=0.05$. Since $R<g<2R$ ($1.01<1.5<2.02$), Corollary~\ref{cor:regimes}(ii) applies, with critical intensity of choice
\begin{equation*}
\beta^{*} = \frac{1}{0.05}\ln\!\left(\frac{1.01}{1.5-1.01}\right) = 20\ln(2.0612) \approx 14.47.
\end{equation*}
Table~\ref{tab:brockhommes} evaluates the equilibrium fraction $n_2^{*}(\beta)$ (Equation~\eqref{eq:nstar}) and the stability eigenvalue $\lambda(\beta)=g\,n_2^{*}(\beta)/R$ (Theorem~\ref{thm:stability}) at several values of $\beta$ spanning this threshold.

\begin{table}[htbp]
\centering
\caption{Equilibrium trend-chaser share and stability eigenvalue as the intensity of choice $\beta$ rises ($R=1.01$, $g=1.5$, $C=0.05$)}
\label{tab:brockhommes}
\begin{tabular}{@{}lrrl@{}}
\toprule
$\beta$ & $n_2^{*}(\beta)$ & $\lambda(\beta)=g\,n_2^{*}/R$ & Fundamental steady state \\
\midrule
$0$ & $0.5000$ & $0.7426$ & stable \\
$5$ & $0.5622$ & $0.8349$ & stable \\
$10$ & $0.6225$ & $0.9246$ & stable \\
$14.47\;(=\beta^{*})$ & $0.6732$ & $1.0000$ & boundary \\
$18$ & $0.7110$ & $1.0559$ & unstable \\
$25$ & $0.7773$ & $1.1545$ & unstable \\
\bottomrule
\end{tabular}
\end{table}

Table~\ref{tab:brockhommes} illustrates the bifurcation directly: the equilibrium trend-chaser share rises smoothly and monotonically from $50\%$ toward (but never reaching) $100\%$ as $\beta$ increases, exactly as Corollary~\ref{cor:regimes}'s proof requires, and the stability eigenvalue crosses $1$ at almost exactly the predicted $\beta^{*}\approx14.47$, confirming $\lambda(\beta^{*})=1$ up to rounding. Below $\beta^{*}$, a small perturbation of the price away from fundamental value damps out; above it, the identical linear mechanism amplifies the same perturbation instead, purely because the equilibrium population has shifted enough toward the trend-following rule to push $g\,n_2^{*}$ past $R$. Nothing about either rule's own forecasting formula, $f_1\equiv0$ or $f_2(x)=1.5x$, changed anywhere in this table.

\section{Discussion: Assumptions, Extensions, and Connections}
\label{sec:discussion}

\textbf{Homogeneous variance.} Assumption (A3) gives every type the same conditional variance $\sigma^2$; this is the standard Brock-Hommes simplification and is what keeps Lemma~\ref{lem:cara}'s demand linear in the forecast $d_{h,t}$ alone. Relaxing it (heterogeneous risk perception across types) is analytically tractable but breaks the clean linear demand and is typically handled numerically in extensions of the base model.

\textbf{Fitness measure.} Equation~\eqref{eq:fitness} uses squared forecast error net of cost, chosen here because it keeps the state space to two lags and yields the closed-form Lemma~\ref{lem:steadystate}. Brock and Hommes's own papers more often use realized trading profit, $(x_{t}-Rx_{t-1})z_{h,t-1}-C_h$, as the fitness measure; this requires one additional lag in the state and complicates the linearization bookkeeping, but does not change the qualitative bifurcation structure of Corollary~\ref{cor:regimes}, since the argument driving that result, both rules tie at the steady state, only the cost breaks the tie, applies to either fitness measure identically.

\textbf{Memory.} A common further extension (also present in Brock and Hommes 1998) replaces Equation~\eqref{eq:fitness} with an exponentially weighted average of current and past fitness, $U_{h,t-1}+\eta U_{h,t-2}^{\text{old}}$ for a memory parameter $\eta\in[0,1)$; larger memory slows how quickly the population reacts to a single period's performance and, as intuition suggests, generally raises the critical $\beta^{*}$ at which instability sets in. This note treats the no-memory case $\eta=0$ throughout.

\textbf{More than two types, and the large type limit.} Section~\ref{sec:twotype} specializes to $H=2$ purely for tractability. Brock, Hommes, and Wagener (2005) study the limit $H\to\infty$ with a continuum of belief types (the ``Large Type Limit''), showing that the discrete switching model converges to a deterministic partial differential equation governing the evolution of the belief distribution, connecting this literature to tools from statistical mechanics.

\textbf{Empirical support.} Boswijk, Hommes, and Manzan (2007) fit a two-type Brock-Hommes-style model directly to more than a century of annual S\&P 500 data using nonlinear least squares, and find a statistically significant, economically large estimated intensity of choice, direct empirical evidence that performance-based switching between a fundamentalist and a trend-following rule is not merely a theoretical curiosity.

\textbf{Relationship to this book's fixed-mix example.} The fundamentalist-chartist model used elsewhere in this book fixes the population shares $(n_f,n_c)$ as free parameters and derives a stability condition on them directly; that model is exactly the special case of Equation~\eqref{eq:deviation} obtained by holding $n_{h,t}$ constant, skipping Sections~\ref{sec:switching}--\ref{sec:stability} of this note entirely. This note supplies the missing piece: it shows that a population split is not merely a modeling choice but can itself be derived as the equilibrium of a rational, evolutionary rule-selection process, and that the resulting equilibrium split is not fixed at all but depends on $\beta$, so that a market can drift from the stable regime of that simpler model into the unstable one purely through gradually more decisive rule-switching, with no change to the underlying trading rules. This is also a candidate mechanism, alongside the endogenous-switching stylized-facts construction used elsewhere in this book, for the same volatility clustering and fat tails Chapter 7's GARCH model captures statistically without explaining.

\textbf{Global dynamics.} Theorem~\ref{thm:stability} and Corollary~\ref{cor:regimes} establish only \emph{local} (linearized) stability of the fundamental steady state; they say nothing about where the system's trajectories go once that steady state loses stability. Brock and Hommes (1997) go considerably further, using the full nonlinear map to show a period-doubling route to chaotic dynamics as $\beta$ increases past $\beta^{*}$, a global nonlinear-dynamics result well beyond a local linearization argument and not reproduced here.

\section{Summary}
\label{sec:summary}

A population of CARA investors with heterogeneous, but commonly-variance, beliefs about a risky asset's payoff clears the market at a price that is a population-weighted average of every belief type's forecast, discounted at the risk-free rate (Section~\ref{sec:model}); recast around a no-bubble fundamental price, this becomes a simple linear recursion in the price's deviation from fundamental, driven by each type's forecasting rule applied to the most recently observed deviation (Section~\ref{sec:fundamental}). Making the population weights themselves an equilibrium object, via a discrete-choice rule that rewards each type in proportion to its recent forecasting accuracy net of cost (Section~\ref{sec:switching}), turns a fixed-mix model into a genuine dynamical system in which prices and population shares co-determine each other. For the canonical two-type case, a costly fundamentalist rule against a free trend-extrapolating rule, the fundamental steady state has a closed-form equilibrium population split (Lemma~\ref{lem:steadystate}) and a closed-form local stability condition (Theorem~\ref{thm:stability}), which together produce a genuine bifurcation in the model's intensity-of-choice parameter $\beta$: for a range of the trend-extrapolation strength $g$, the market is stable at low $\beta$ and loses stability entirely once $\beta$ crosses an explicit critical value $\beta^{*}$ (Corollary~\ref{cor:regimes}), purely because a more performance-sensitive population rationally migrates toward the cheaper of two equally accurate rules. This is the local mechanism underlying Brock and Hommes's (1997) ``rational route to randomness,'' in which fully rational, evolutionary rule selection can generate market dynamics that are statistically indistinguishable from noise.

\section*{References}
\begin{itemize}
\item Brock, W. A., and Hommes, C. H. (1997). ``A Rational Route to Randomness.'' \textit{Econometrica}, 65(5), 1059--1095.
\item Brock, W. A., and Hommes, C. H. (1998). ``Heterogeneous Beliefs and Routes to Chaos in a Simple Asset Pricing Model.'' \textit{Journal of Economic Dynamics and Control}, 22(8-9), 1235--1274.
\item Brock, W. A., Hommes, C. H., and Wagener, F. O. O. (2005). ``Evolutionary Dynamics in Markets with Many Trader Types.'' \textit{Journal of Mathematical Economics}, 41(1-2), 7--42.
\item Boswijk, H. P., Hommes, C. H., and Manzan, S. (2007). ``Behavioral Heterogeneity in Stock Prices.'' \textit{Journal of Economic Dynamics and Control}, 31(6), 1938--1970.
\item Hommes, C. H. (2006). ``Heterogeneous Agent Models in Economics and Finance.'' In \textit{Handbook of Computational Economics}, Vol.\ 2, edited by L.\ Tesfatsion and K.\ L.\ Judd, 1109--1186. Elsevier.
\item Lux, T. (1995). ``Herd Behaviour, Bubbles and Crashes.'' \textit{Economic Journal}, 105(431), 881--896.
\end{itemize}

\end{document}
